\documentclass[11pt,a4paper]{article}

\usepackage[T1]{fontenc}
\usepackage[utf8]{inputenc}
\usepackage{lmodern}
\usepackage[a4paper,top=2.4cm,bottom=2.2cm,left=2.2cm,right=2.2cm]{geometry}
\usepackage{amsmath}
\usepackage{amssymb}
\usepackage{graphicx}
\graphicspath{{./}}
\usepackage{array}
\usepackage{makecell}
\renewcommand{\arraystretch}{1.25}
\usepackage{enumitem}
\usepackage{fancyhdr}
\usepackage{xcolor}

\newcommand{\ohms}{\ensuremath{\Omega}}
\newcommand{\degC}{\ensuremath{^{\circ}}C}
\newcommand{\dg}{\ensuremath{^{\circ}}}
\newcommand{\pow}[1]{\ensuremath{^{#1}}}

\newcommand{\topiclabel}[1]{{\hfill\normalfont\small\itshape\textcolor{black!55}{#1}}}

\newlength{\anslineskip}
\setlength{\anslineskip}{8mm}

\newcounter{ansln}
\newcommand{\Alines}[2]{%
  \par\vspace{1mm}%
  \setcounter{ansln}{1}%
  \loop\ifnum\value{ansln}<#1\relax
    \noindent\mbox{}\dotfill\par\vspace{\anslineskip}%
    \stepcounter{ansln}%
  \repeat
  \noindent\mbox{}\dotfill\hspace{0.6em}\mbox{#2}\par
  \vspace{1mm}%
}

\newcommand{\plainline}{\par\noindent\mbox{}\dotfill\par\vspace{\anslineskip}}

\newcommand{\labelline}[1]{\par\vspace{1mm}\noindent #1\hspace{0.5em}\dotfill\par\vspace{\anslineskip}}

\newcommand{\ansval}[3]{%
  \par\vspace{5mm}%
  \noindent\hspace*{3.5cm}#1\hspace{0.5em}\dotfill\hspace{0.5em}#2\hspace{0.5em}\makebox[1.1cm][r]{#3}\par
  \vspace{3mm}%
}

\renewcommand{\markright}[1]{\par\nopagebreak\noindent\hspace*{\fill}\mbox{#1}\par\vspace{1mm}}

% Per-question head: \examq{paper-id}{number}{topics}{marks}
% Renders: paper-id - Q<number> - topics - marks  (one line, same font).
% Same command and same rendered line as the answer paper.
\newcommand{\swmarks}{}
\newcommand{\examq}[4]{%
  \gdef\swmarks{#4}%
  \par\addvspace{16pt}\noindent
  {\large\bfseries #1 - Q#2 - #3 - #4}\par\vspace{5pt}%
}
% \total now takes NO argument: it prints the marks stored by \examq.
\newcommand{\total}{\par\vspace{2mm}\noindent\hspace*{\fill}[Total:\ \swmarks]\par\vspace{2mm}}

\newcommand{\qfig}[3][0.5\linewidth]{%
  \par\vspace{2mm}
  \begin{center}%
    \IfFileExists{#2}%
      {\includegraphics[width=#1]{#2}}%
      {\fbox{\parbox[c][26mm][c]{#1}{\centering\ttfamily\footnotesize #2}}}\\[4pt]
    \textbf{#3}%
  \end{center}%
  \vspace{2mm}\par
}

\newlist{parts}{enumerate}{1}
\setlist[parts]{label=\textbf{(\alph*)},leftmargin=2.1em,labelsep=0.6em,itemsep=10pt,topsep=6pt,parsep=4pt}

\newlist{subparts}{enumerate}{1}
\setlist[subparts]{label=\textbf{(\roman*)},leftmargin=2.3em,labelsep=0.6em,itemsep=10pt,topsep=6pt,parsep=4pt}

\pagestyle{fancy}
\fancyhf{}
\fancyhead[L]{\footnotesize 4037/21/M/J/25}
\fancyhead[C]{\footnotesize Cambridge O Level Additional Mathematics -- Paper 2}
\fancyhead[R]{\footnotesize May/June 2025}
\fancyfoot[C]{\footnotesize\thepage}
\renewcommand{\headrulewidth}{0.4pt}
\renewcommand{\footrulewidth}{0pt}

\setlength{\parindent}{0pt}

\begin{document}

%==================== QUESTION 1 ====================
\examq{4037/21/M/J/25}{1}{Quadratic Functions}{5}
A curve has equation $y=x^{2}+2x-3$.
\begin{parts}
  \item Use the method of completing the square to find the coordinates of the stationary point on the curve.
  \markright{[3]}
  \vspace{5cm}

  \item On the axes, sketch the curve, stating the intercepts with the coordinate axes.
  \markright{[2]}
  \qfig[0.6\linewidth]{4037_21_M_J_25_fig1.png}{}
\end{parts}
\total

%==================== QUESTION 2 ====================
\examq{4037/21/M/J/25}{2}{Quadratic Functions}{5}
In this question, $k$ is a constant.\\
It is given that $2x^{2}+(3k-2)x+k=0$ has roots that are real and distinct.

Find the set of possible values of $k$.
\markright{[5]}
\vspace{10cm}
\total

%==================== QUESTION 3 ====================
\examq{4037/21/M/J/25}{3}{Permutations \& Combinations}{5}
A sports club has the following members.
\begin{center}
5 runners, 4 swimmers, 3 gymnasts
\end{center}
\begin{parts}
  \item These members stand in a straight line.

  Find the number of ways that this can be done when all the runners stand together, all the swimmers stand together and all the gymnasts stand together.
  \markright{[2]}
  \vspace{4cm}

  \item Four of these members are selected for an event.

  Find the number of ways that this can be done when at least one runner, at least one swimmer and at least one gymnast must be selected.
  \markright{[3]}
  \vspace{5cm}
\end{parts}
\total

%==================== QUESTION 4 ====================
\examq{4037/21/M/J/25}{4}{Straight-Line Graphs}{7}
The coordinates of points $A$, $B$, $C$ and $D$ are as follows.
\begin{center}
$A(-4,\,3)$\quad $B(6,\,-9)$\quad $C(15,\,10)$\quad $D(14,\,-1)$
\end{center}
The line $L$ has equation $y=11x-75$.\\
The perpendicular bisector of the line $AB$ meets $L$ at the point $E$.

Find the area of triangle $CDE$.
\markright{[7]}
\vspace{17cm}
\total

%==================== QUESTION 5 ====================
\examq{4037/21/M/J/25}{5}{Differentiation}{4}
Given that $y=x^{2}\tan\dfrac{x}{2}$, use calculus to find the approximate change in $y$ as $x$ increases from $\dfrac{\pi}{3}$ to $\dfrac{\pi}{3}+h$, where $h$ is small.
\markright{[4]}
\vspace{18cm}
\total

%==================== QUESTION 6 ====================
\examq{4037/21/M/J/25}{6}{Binomial Theorem}{6}
\begin{parts}
  \item Using an appropriate quadratic factorisation, find the first three terms in the binomial expansion of $\left(9x^{2}+12x+4\right)^{5}$, in ascending powers of $x$.\\
  You must simplify your coefficients.
  \markright{[4]}
  \vspace{9cm}

  \item Find the term independent of $x$ in the expansion of $\left(\dfrac{6}{x^{2}}+\dfrac{x^{4}}{2}\right)^{12}$.
  \markright{[2]}
  \vspace{8cm}
\end{parts}
\total

%==================== QUESTION 7 ====================
\examq{4037/21/M/J/25}{7}{Functions \textperiodcentered\ Logarithmic \& Exponential Functions}{8}
\begin{parts}
  \item The function f is defined by $\mathrm{f}(x)=2\mathrm{e}^{-x}+3$ for $x\in\mathbb{R}$.

  On the axes, sketch the graph of $y=\mathrm{f}(x)$ and hence, on the same axes, sketch the graph of $y=\mathrm{f}^{-1}(x)$ .\\
  Show clearly
  \begin{itemize}[leftmargin=2em,itemsep=2pt,topsep=2pt]
    \item the positions of any points where your graphs meet the coordinate axes
    \item the positions of any asymptotes.
  \end{itemize}
  \markright{[4]}
  \qfig[0.6\linewidth]{4037_21_M_J_25_fig2.png}{}

  \item The function g is defined by $\mathrm{g}(x)=2-\dfrac{3}{\mathrm{e}^{x}+2}$ for $x\geqslant 0$.

  Given that $\mathrm{g}^{-1}$ exists, find an expression for $\mathrm{g}^{-1}(x)$ and state its domain.
  \markright{[4]}
  \vspace{17cm}
\end{parts}
\total

%==================== QUESTION 8 ====================
\examq{4037/21/M/J/25}{8}{Trigonometry}{9}
\begin{parts}
  \item Solve the equation $(2-3\cot x)\cos x=0$ for $0<x\leqslant\dfrac{\pi}{2}$.
  \markright{[3]}
  \vspace{18cm}

  \item Solve the equation $2\operatorname{cosec}(2\theta+1)-12\sin(2\theta+1)=5$, where $\theta$ is in radians and $-1\leqslant\theta\leqslant 2$.
  \markright{[6]}
  \vspace{18cm}
\end{parts}
\total

%==================== QUESTION 9 ====================
\examq{4037/21/M/J/25}{9}{Circular Measure \textperiodcentered\ Differentiation}{12}
In this question, all lengths are in centimetres and all angles are in radians.
\qfig[0.85\linewidth]{4037_21_M_J_25_fig3.png}{}
The diagram shows a sector $AOB$ of a circle, centre $O$, radius 15.

Angle $AOB=\dfrac{6\pi}{5}$.

The sector is made into a cone with points $A$ and $B$ touching, as shown.
\begin{parts}
  \item Find the curved surface area of the cone.
  \markright{[2]}
  \vspace{3cm}
\end{parts}

The top of the cone is a horizontal circle.

\begin{parts}[resume]
  \item Find the circumference of the circular top.
  \markright{[2]}
  \vspace{3cm}

  \item Hence find the radius of the circular top and the perpendicular height of the cone.
  \markright{[2]}
  \vspace{5cm}

  \item Water is poured into the cone.\\
  When the depth of the water in the cone is $h$, the radius of the circular top of the water is $r$.
  \begin{subparts}
    \item Find an expression for $r$ in terms of $h$.
    \markright{[1]}
    \vspace{4cm}

    \item The water is poured into the cone at a constant rate of $27\,\mathrm{cm}^{3}$ per second.

    Find the rate at which the depth of the water is rising when the depth of the water is 4.
    \markright{[5]}
    \vspace{10cm}
  \end{subparts}
\end{parts}
\total

%==================== QUESTION 10 ====================
\examq{4037/21/M/J/25}{10}{Calculus for Kinematics}{11}
A particle $P$ moves in a straight line.\\
$t$ seconds after passing a fixed point, $O$, the acceleration of $P$, $a\,\mathrm{ms}^{-2}$, is given by
\[
\begin{aligned}
  a&=t^{2}-2 &&\text{for } 0\leqslant t\leqslant 4\\
  a&=19-5\mathrm{e}^{8-2t} &&\text{for } 4\leqslant t\leqslant 10.
\end{aligned}
\]
When $t=3$, the velocity of $P$ is $-\dfrac{1}{3}\,\mathrm{ms}^{-1}$ and its displacement from $O$ is $-\dfrac{1}{4}\,\mathrm{m}$.
\begin{parts}
  \item \begin{subparts}
    \item Find the velocity of $P$ when $t=4$.
    \markright{[3]}
    \vspace{6cm}

    \item Find the displacement of $P$ from $O$ when $t=4$.
    \markright{[3]}
    \vspace{9cm}
  \end{subparts}

  \item Find the displacement of $P$ from $O$ when $t=10$.
  \markright{[5]}
  \vspace{17cm}
\end{parts}
\total

\begin{center}\textbf{Question 11 is printed on the next page.}\end{center}

%==================== QUESTION 11 ====================
\examq{4037/21/M/J/25}{11}{Arithmetic \& Geometric Progressions}{8}
A geometric progression has first term $a$ and common ratio $r$, where $r>0$.\\
The sum of the 2nd and 3rd terms of the progression is 168.\\
The sum of the 4th and 5th terms of the progression is 94.5.
\begin{parts}
  \item Find the 6th term of the progression.
  \markright{[7]}
  \vspace{13cm}

  \item Find the sum to infinity of the progression.
  \markright{[1]}
  \vspace{5cm}
\end{parts}
\total

\end{document}